\lecture{4.5}{Comparing G-Computation and IPTW}
\label{lec:4-5}

\subsection{When Empirical G-Computation and IPTW Coincide}

In the fully discrete setting, if both nuisance functions are estimated by empirical
subgroup proportions and means, empirical g-computation and empirical IPTW give the
same point estimate. This equality is not a general property of all g-computation and
IPTW estimators. It is a special algebraic feature of saturated discrete estimation.

\begin{proposition}[Equivalence under empirical discrete estimation]
Suppose $W$ is discrete, every relevant $(A=a,W=w)$ cell is nonempty, and
\[
\Qbar_n(a,w)
=
\frac{\sum_{i=1}^nY_i\ind(A_i=a,W_i=w)}
       {\sum_{i=1}^n\ind(A_i=a,W_i=w)},
\qquad
g_{n,a}(w)
=
\frac{\sum_{i=1}^n\ind(A_i=a,W_i=w)}
       {\sum_{i=1}^n\ind(W_i=w)}.
\]
Then
\[
\frac{1}{n}\sum_{i=1}^n \Qbar_n(a,W_i)
=
\frac{1}{n}\sum_{i=1}^n
\frac{\ind(A_i=a)}{g_{n,a}(W_i)}Y_i.
\]
\end{proposition}

\begin{proof}
Let $n_w=\sum_i\ind(W_i=w)$ and $n_{aw}=\sum_i\ind(A_i=a,W_i=w)$.
The g-computation estimator can be written as
\[
\sum_w \Qbar_n(a,w)\frac{n_w}{n}.
\]
Substituting the empirical subgroup mean,
\[
\sum_w
\left\{
\frac{\sum_iY_i\ind(A_i=a,W_i=w)}{n_{aw}}
\right\}
\frac{n_w}{n}.
\]
Substituting $g_{n,a}(w)=n_{aw}/n_w$ gives
\[
\frac{1}{n}\sum_w
\sum_i
Y_i\ind(A_i=a,W_i=w)
\frac{1}{g_{n,a}(w)}.
\]
For each observation $i$, only the term with $w=W_i$ contributes, so the expression
equals
\[
\frac{1}{n}\sum_{i=1}^n
\frac{\ind(A_i=a)}{g_{n,a}(W_i)}Y_i,
\]
which is the IPTW estimator.
\end{proof}

\subsection{Why the Equivalence Is Special}

The proof above uses two facts. First, $\Qbar_n(a,w)$ is the empirical mean in exactly
the same strata used to estimate $g_{n,a}(w)$. Second, the empirical covariate
distribution gives each stratum weight $n_w/n$. If either nuisance function is
estimated by a model that smooths across strata, the exact cancellation in the proof
usually disappears.

\begin{keyidea}
G-computation and IPTW are equivalent as identification formulas under the causal
assumptions. Their finite-sample estimators need not be equal, because the fitted
outcome regression and fitted propensity score can approximate different parts of the
observed-data distribution in different ways.
\end{keyidea}

\begin{figure}[H]
\centering
\begin{tikzpicture}[
  node distance=1.45cm,
  box/.style={draw, rounded corners=2pt, minimum width=3.2cm, minimum height=.8cm},
  arr/.style={-{Latex[length=2.5mm]}, thick}
]
\node[box] (id) {Same identified\\ $\psi(a)$};
\node[box, below left=1.1cm and 1.2cm of id] (gcomp) {Estimate\\ $\Qbar(a,w)$};
\node[box, below right=1.1cm and 1.2cm of id] (iptw) {Estimate\\ $g_a(w)$};
\node[box, below=1.2cm of gcomp] (pg) {Average\\ predictions};
\node[box, below=1.2cm of iptw] (pi) {Average\\ weighted outcomes};
\draw[arr] (id) -- (gcomp);
\draw[arr] (id) -- (iptw);
\draw[arr] (gcomp) -- (pg);
\draw[arr] (iptw) -- (pi);
\draw[arr, dashed] (pg) -- node[below] {equal only in special cases} (pi);
\end{tikzpicture}
\caption{The two methods estimate the same causal parameter through different nuisance functions.}
\label{fig:ch4-gcomp-iptw-comparison}
\end{figure}

\subsection{High-Dimensional Discrete Covariates}

Empirical stratification becomes fragile when $W$ has many discrete components.
If $W=(W_1,\ldots,W_p)$ and each $W_j$ is binary, then $W$ can take $2^p$ possible
values. Even moderate $p$ can create many sparse or empty strata.

Sparse strata create two problems. For g-computation, $\Qbar_n(a,w)$ is unstable when
only a few observations have $A=a$ and $W=w$, and undefined when no such observation
exists. For IPTW, $g_{n,a}(w)$ is unstable when treatment counts are small, and the
weights $1/g_{n,a}(w)$ can become extremely large.

\subsection{Continuous Covariates}

If $W$ contains a continuous component, exact empirical stratification usually fails.
In a continuous distribution, the probability that two individuals have exactly the
same covariate vector is often zero. Then the subgroup
\[
\{i:W_i=w\}
\]
may contain only one observation or none, making empirical subgroup means unusable.

This is not a conceptual problem for the estimand. The counterfactual mean
\[
\psi(a)=\E\{\Qbar(a,W)\}
\]
is still well-defined. The difficulty is statistical: we must estimate the smooth
conditional mean or treatment probability from finite data.

\subsection{Need for Regression Smoothing}

Regression smoothing borrows information across similar covariate values. Instead of
estimating a separate mean for every exact $w$, a model might assume
\[
\Qbar(a,w)=\beta_0+\beta_1a+\beta_2^\top w
\]
or
\[
g_1(w)=\expit(\alpha_0+\alpha_1^\top w).
\]
These restrictions reduce variance, but they introduce the possibility of model
misspecification. A simple main-terms model may be too rigid if the true outcome
regression has nonlinearities or treatment-covariate interactions.

The central estimation tradeoff is therefore:
\[
\text{less smoothing} \Rightarrow \text{lower modeling bias but higher variance},
\]
\[
\text{more smoothing} \Rightarrow \text{lower variance but higher risk of bias}.
\]

\subsection{Variance and Instability}

IPTW estimators are especially sensitive to near-violations of positivity. If
$g_{n,1}(W_i)=0.02$, then a treated individual receives weight $50$. A small number of
such observations can strongly influence the weighted mean.

G-computation can also be unstable or biased, but the failure mode is different. A
poor outcome regression may extrapolate into covariate-treatment combinations with
little support. For example, if very old individuals are almost never treated, then
$\Qbar_n(1,w)$ for very old $w$ is largely determined by model assumptions rather
than direct data.

\begin{estimationbox}
Practical comparison:
\begin{itemize}
\item Empirical g-computation and empirical IPTW coincide in low-dimensional
discrete settings with nonempty cells.
\item Regression-based g-computation and regression-based IPTW usually differ.
\item G-computation relies most directly on the quality of the outcome regression.
\item IPTW relies most directly on the quality of the propensity score and on stable
positivity.
\item Large weights, empty cells, and extrapolation are diagnostic signs that the
data may not contain enough information for the desired causal contrast without
strong modeling assumptions.
\end{itemize}
\end{estimationbox}
